\section{Data, versions and preprocessing}\label{sec:data}

\subsection{Raw inputs}
Table~\ref{tab:datasources} lists every raw dataset with its version, size and file date. The manifest
\file{docs/DATA\_MANIFEST.csv} records each of the \manifestFiles\ raw files (\manifestGB\,GB) with its SHA-256
checksum, so any later run can confirm that it uses exactly the same inputs (Section~\ref{sec:repro}). The files
come from two snapshots:
\begin{itemize}
\item the mid-evaluation snapshot (5 October 2026): climate, SPAM, the original calendars, FAOSTAT, USDA and
Natural Earth;
\item an end-evaluation snapshot (17 November 2026): the calendars of the new crops, the WorldClim normals and
CMIP6 files, the FAOSTAT Supply Utilization Accounts, the admin-1 boundaries and the Indian survey files.
\end{itemize}
FAOSTAT bulk files are revised over time, so the numbers are exactly reproducible only from these archived files.

{\scriptsize\begin{longtable}{L{0.275}L{0.131}L{0.137}L{0.191}rrL{0.096}}
\caption[Raw inputs]{Raw inputs (from \texttt{docs/DATA\_MANIFEST.csv}, which lists every file with its SHA-256 checksum).}\label{tab:datasources}\\
\toprule
Dataset & Version & Source & Used for & Files & MB & File date \\
\midrule\endfirsthead
\toprule
Dataset & Version & Source & Used for & Files & MB & File date \\
\midrule\endhead
\bottomrule\endfoot
WorldClim 2.1 monthly weather (CRU-TS 4.09 downscaled), 10$'$; elevation; 1970--2000 precipitation normals & CRU-TS 4.09; WorldClim 2.1 & \citet{fick2017,harris2020} & monthly precipitation 1970--2024; elevation & 8 & 1{,}209 & 2026-10-05 \\
WorldClim 2.1 monthly maximum temperature, 10$'$ & CRU-TS 4.09 & \citet{fick2017,harris2020} & $T_{\max}$ 1970--2024 & 6 & 748 & 2026-10-05 \\
WorldClim 2.1 monthly minimum temperature, 10$'$ & CRU-TS 4.09 & \citet{fick2017,harris2020} & $T_{\min}$ 1970--2024 & 6 & 725 & 2026-10-05 \\
WorldClim 2.1 monthly normals $T_{\max}$, $T_{\min}$, 10$'$ & 1970--2000 & \citet{fick2017} & baseline of the change factors & 2 & 75 & 2026-11-17 \\
WorldClim 2.1 CMIP6 downscaled climatology, MRI-ESM2-0, SSP2-4.5, 10$'$ & 2041--2060 & \citet{wcmip6} & future change factors & 3 & 54 & 2026-11-17 \\
SPAM 2010 harvested area (all, irrigated, rainfed), 5$'$ & v2.0 & \citet{spam2010} & crop area weights, rainfed share & 2 & 144 & 2026-10-05 \\
Crop calendar dataset (planting, harvest) & Sacks et al.\ 2010 & \citet{sacks2010} & growing season and flowering window & 20 & 70 & 2026-10-05 / 2026-11-17 \\
FAOSTAT bulk: crop production (QCL); Food Balances 2010-- and --2013; Food Security indicators & bulk files of 5 Oct 2026 & \citet{faostat} & panels, yields, nutrient supply, energy requirement & 4 & 163 & 2026-10-05 \\
FAOSTAT Supply Utilization Accounts, crops & bulk file of 17 Nov 2026 & \citet{faosua} & chickpea and lentil supply & 1 & 89 & 2026-11-17 \\
USDA FoodData Central, SR Legacy & April 2018 release & \citet{usda2019} & zinc and iron per 100 g & 1 & 6 & 2026-10-05 \\
Natural Earth admin-0 countries, 1:50m & -- & Natural Earth (public domain) & country masks, continents & 1 & 1 & 2026-10-05 \\
Natural Earth admin-1 states/provinces, 1:10m & -- & Natural Earth (public domain) & Telangana mask & 1 & 15 & 2026-11-17 \\
K\"oppen--Geiger 1991--2020, 0.5\textdegree & Beck et al.\ 2023 & \citet{beck2018} & check of our K\"oppen groups & 1 & 0 & 2026-10-05 \\
Official India/Telangana/Hyderabad files: MoSPI HCES 2023--24, NSS 592 and 594, NFHS-5, CNNS, ICMR-NIN IFCT 2017, Telangana DES & see Appendix~\ref{app:indiasources} & Section~\ref{sec:india} & Hyderabad and Telangana analysis & 18 & 99 & 2026-11-17 \\
\end{longtable}}


\subsection{Gridded climate, crop area and calendars}
All gridded work is done on a 10 arc-minute global grid (about 18\,km at the equator).
\begin{itemize}
\item \textbf{Climate.} WorldClim~2.1 historical monthly weather (CRU-TS 4.09 downscaled) gives monthly
$T_{\max}$, $T_{\min}$ and precipitation $P$ for 1970--2024; $T_{\text{mean}} = (T_{\max}+T_{\min})/2$. Monthly data
cannot show individual hot days, a limitation discussed in Section~\ref{sec:limits}. In cell-months with no nearby
station, CRU-TS falls back to climatology, which weakens year-to-year rainfall variation in data-poor regions.
\item \textbf{Crop area.} SPAM~2010 harvested area (all, irrigated, rainfed; 5$'$) is summed to the 10$'$ grid. It
gives each cell's crop hectares $A_i$ and rainfed hectares $R_i$. Cells are assigned to countries with the
Natural Earth admin-0 boundaries.
\item \textbf{Calendars.} The Sacks et al.\ dataset gives planting and harvest days per cell; observed values are
used first and gap-filled values only as a fallback. Chickpea, lentil and rapeseed-mustard use a
\textbf{cool-season rule}, for three reasons:
  \begin{itemize}
  \item Sacks et al.\ have no chickpea or lentil calendar;
  \item their generic ``Pulses'' calendar is 95\% gap-filled in India;
  \item that calendar places these winter (rabi) crops in the monsoon season.
  \end{itemize}
  The rule takes the crop's own observed calendar where one exists (winter rapeseed in China, India and Europe),
  then the spring-wheat calendar where winters are cold (K\"oppen group D), then the wheat calendar elsewhere. This
  gives rabi sowing in South Asia, autumn sowing in Australia and spring sowing in Canada and Russia. The plain
  Sacks calendar is kept as a robustness test.
\item \textbf{Elevation and climate zones.} Elevation and K\"oppen groups (computed from WorldClim with the
published rules, \citealp{beck2018}) only \emph{describe} the panel countries; they never select them.
\end{itemize}

\subsection{Food balance sheets, supply utilization accounts and composition}\label{sec:fbsmap}
\paragraph{Nutrient supply.} National supply of energy, protein and fat per person per day comes from the FAOSTAT
food balance sheets (FBS). The new series (2010--2023) is used for importance and adequacy. The old series
(1961--2013) is used only for the historical pass-through regression. The two series use different methods, so
they are analysed separately and never joined.

Zinc and iron are not in the FBS. They are computed item by item as FBS food quantity multiplied by a composition
value from USDA FoodData Central (SR Legacy).
\begin{itemize}
\item Each FBS item is mapped to the raw, unprocessed USDA food closest to its primary-commodity equivalent
(\file{code/fbs\_usda\_map.csv}). This overstates zinc and iron for cereals eaten refined and for meat and fish in
carcass or live weight.
\item Five minor items with no reasonable analogue and negligible supply are set to zero.
\item The sum of our item-level energy equals FAO's Grand Total to four decimals, so no item is missing.
\end{itemize}

\paragraph{Chickpea and lentil.} The FBS fold chickpea and lentil into ``Pulses, Other and products''. Their
supply comes instead from the FAOSTAT Supply Utilization Accounts (SUA), the tables from which the FBS are
compiled (items ``Chick peas, dry'' and ``Lentils, dry'').
\begin{itemize}
\item SUA domestic supply is production $+$ imports $-$ exports $-$ stock variation. We checked this identity
against the sum of all uses in the raw SUA file: it holds in 100\% of chickpea and lentil country-years, while the
opposite sign holds in only half.
\item The SUA rows give each crop's contribution. They are \emph{never} added to national totals, which already
contain them inside the FBS aggregate.
\item They are not entered as extra items in the food-supply diversity calculation.
\item Across country-years the SUA chickpea and lentil energy stays within the parent FBS aggregate in all but
0.1\% of cases (Table~\ref{tab:checks}).
\end{itemize}
How each crop maps to FBS and SUA items is in Table~\ref{tab:cropmap}.

\paragraph{Requirements.} Requirements combine FAOSTAT's country- and year-specific Average Dietary Energy
Requirement (ADER) with the WHO/FAO recommended intakes (Section~\ref{sec:adequacy}).

{\scriptsize\begin{longtable}{L{0.105}L{0.072}L{0.124}L{0.059}L{0.137}L{0.170}L{0.098}L{0.235}}
\caption[How each crop is identified in each dataset]{How each crop is identified in each dataset (from \texttt{code/config.py} and \texttt{code/fbs\_usda\_map.csv}).}\label{tab:cropmap}\\
\toprule
Crop & Group & FAOSTAT item & SPAM & Sacks calendar & FBS items & SUA item & USDA composition (Zn, Fe) \\
\midrule\endfirsthead
\toprule
Crop & Group & FAOSTAT item & SPAM & Sacks calendar & FBS items & SUA item & USDA composition (Zn, Fe) \\
\midrule\endhead
\bottomrule\endfoot
Rice & cereal & Rice & RICE & Rice & Rice and products; Rice (Milled Equivalent) & -- & Rice, white, long-grain, regular, raw, unenriched \\
Wheat & cereal & Wheat & WHEA & Wheat.Winter, Wheat & Wheat and products & -- & Wheat, hard red winter \\
Maize & cereal & Maize (corn) & MAIZ & Maize & Maize and products & -- & Corn grain, yellow \\
Soybean & oilseed & Soya beans & SOYB & Soybeans & Soyabeans; Soyabean Oil & -- & Soybeans, mature seeds, raw; Oil, soybean, salad or cooking \\
Chickpea & pulse & Chick peas, dry & CHIC & Pulses (cool-season rule) & -- & Chick peas, dry & Chickpeas (garbanzo beans, bengal gram), mature seeds, raw \\
Lentil & pulse & Lentils, dry & LENT & Pulses (cool-season rule) & -- & Lentils, dry & Lentils, raw \\
Common bean & pulse & Beans, dry & BEAN & Pulses & Beans & -- & Beans, kidney, all types, mature seeds, raw \\
Groundnut & oilseed & Groundnuts, excluding shelled & GROU & Groundnuts & Groundnuts; Groundnut Oil & -- & Peanuts, all types, raw; Oil, peanut, salad or cooking \\
Rapeseed-mustard & oilseed & Rape or colza seed & RAPE & Rapeseed.Winter (cool-season rule) & Rape and Mustardseed; Rape and Mustard Oil & -- & Spices, mustard seed, ground; Oil, canola \\
\end{longtable}}


\subsection{Future climate}
The future input is the WorldClim~2.1 downscaled CMIP6 climatology for MRI-ESM2-0, SSP2-4.5, 2041--2060. These are
20-year mean monthly values of $T_{\max}$, $T_{\min}$ and $P$. Their baseline is the WorldClim 1970--2000 monthly
normals. How they are combined with the observed years is described in Section~\ref{sec:periods}. Only one model
and one scenario are used, so the future results show direction under that model, not the range of possible
futures.

\subsection{India, Telangana and Hyderabad}
The local analysis uses official files listed in Appendix~\ref{app:indiasources}:
\begin{itemize}
\item MoSPI Household Consumption Expenditure Survey 2023--24 (state $\times$ rural/urban item quantities);
\item NSS reports 592 and 594 (expenditure and nutrient intake by state and sector);
\item NFHS-5 2019--21 (India, Telangana and the Hyderabad district fact sheet);
\item CNNS 2016--18 (Telangana micronutrient status);
\item the ICMR-NIN Indian Food Composition Tables 2017;
\item the Telangana Directorate of Economics and Statistics crop statistics.
\end{itemize}
Each number taken from these files is stored once, with its table or page, in
\file{data/external/india\_hyderabad\_indicators.csv}.
